\begin{problem}[Broken IND-CPA to IND-CCA transformation]\label{p1}
    \leavevmode\par
    \begin{definition*}
        A MAC system is a pair of PPT algorithms $\Sign$, a signing algorithm, and $\Vrfy$, a verification algorithm.
        \begin{enumerate}[label=(\alph*)]
            \item $\Sign$ takes as input a key $k$ and a message $m$ and output a tag $t$.
            \item $\Vrfy$ is a deterministic algorithm that takes as input a key $k$, a message $m$ and a tag $t$ and outputs $\accept$ or $\reject$.
        \end{enumerate}
        We require that tags generated by $\Sign$ are always accepted by $\Vrfy$; that is, the MAC must satisfy the following correctness property: for all keys $k$ and for all messages $m$,
        \[
            \Pr[\Vrfy_k(m, \Sign_k(m))=\accept]=1.
        \]
        As usual, We denote by $\Kcal$ the key space, $\Mcal$ the message space and $\Tcal$ the tag space.
    \end{definition*}

    \begin{definition*}[MAC security with verification queries]
        For a given MAC system $\Ical = (\Sign, \Vrfy)$, the MAC experiment runs as follows:
        \begin{enumerate}[label=(\alph*)]
            \item The challenger picks a random $k\Usample \Kcal$.
            \item $\Adv$ queries the challenger several times. For $i = 1, 2,\ldots $, the $i$th signing query is a message  $m_i\in\Mcal$.
            \item Given $m_i$ the challenger computes a tag $t_i\sample \Sign_k(m_i)$ and returns $t_i$ to $\Adv$.
            \item Eventually $\Adv$ outputs a candidate forgery pair $(m, t)\in \Mcal\times \Tcal$ that is not among the signed pairs, i.e.,
            \[
                (m,t)\notin \{(m_1, t_1), (m_2, t_2),\ldots\}.
            \]
            \item $\Adv$ wins iff the challenger ever responds to a verification query with $\accept$, i.e., $\Vrfy_k(m,t)=\accept$.
        \end{enumerate}
    \end{definition*}

    \begin{definition*}
        We say that a MAC system $\Ical$ is secure if for all PPT adversary $\Adv$, there is a negligible function $\mu$ such that in the above experiment MAC
        \[
            \Pr[\Adv \textrm{ wins}] \le \mu(n).
        \]
    \end{definition*}

    Consider the following attempt at transformation a CPA-secure scheme to an IND-CCA-secure one.
    Let $(\Gen, \Enc, \Dec)$ be an IND-CPA-secure PKE defined over $\Kcal\times \Mcal, \Ccal$ and let $(\Sign, \Vrfy)$ be a secure MAC with key space $\Kcal$. We construct a new PKE scheme $\PKE' = (\Gen', \Enc', \Dec')$ with message space $\Mcal$, as follows:
    \begin{simplebox}
        \small
        \noindent
        \begin{minipage}[t]{0.3\textwidth}
            \begin{algorithmic}[1]
                \Statex \textbf{$\Gen'(1^n)$}
                \Indent
                \State \Return $\Gen(1^n)$
                \EndIndent
            \end{algorithmic}
        \end{minipage}
        \hfill
        \begin{minipage}[t]{0.3\textwidth}
            \begin{algorithmic}[1]
                \Statex \textbf{$\Enc_\pk'(m)$}
                \Indent
                    \State $k\Usample \Kcal$
                    \State $ct\sample \Enc_\pk(k,m)$
                    \State $t\sample \Sign_k(ct)$
                    \State \Return $ct' = (ct, t)$
                \EndIndent
            \end{algorithmic}
        \end{minipage}
        \hfill
        \begin{minipage}[t]{0.38\textwidth}
            \begin{algorithmic}[1]
                \Statex \textbf{$\Dec_\sk'(ct'=(ct, t))$}
                \Indent
                    \State $(k,m)\coloneqq \Dec_\sk (ct)$
                    \If{$\Vrfy_k(ct, t)=\accept$}
                        \State \Return $m$
                    \Else 
                        \State \Return $\bot$
                    \EndIf
                \EndIndent
            \end{algorithmic}
        \end{minipage}
    \end{simplebox}

    One might expect this scheme to be IND-CCA-secure because a change to a ciphertext  $(ct, t)$ will invalidate the MAC tag $t$. Show that this is INCORRECT. i,e,. show an IND-CPA-secure PKE scheme $\Pi = (\Gen, \Enc, \Dec)$ for which $\Pi' = (\Gen', \Enc', \Dec')$ is not IND-CCA-secure (for any choice of MAC).
\end{problem}



\begin{answer}
    To show this, we provide a counterexample.\\
    Let $\Pi^0=(\Gen^0, \Enc^0, \Dec^0)$ be a $\delta$-correct, $\gamma$-spread, IND-CPA-secure PKE scheme over $(\Kcal\times \Mcal, \Ccal)$, where $\delta$ and $2^{-\gamma}$ are negligible. Also, let  $\Ical=(\Sign, \Vrfy)$ be a secure MAC system with message space $\Ccal \times \Ccal$, key space $\Kcal$, and tag space $\Tcal$.

    Fix $k_n^0\in\Kcal_n$ and $m_n^0\in\Mcal_n$ for each $n\in\Nbb$. We construct an IND-CPA-secure PKE scheme $\Pi=(\Gen, \Enc, \Dec)$ over $(\Kcal\times \Mcal, \Ccal\times \Ccal)$ as follows:
    \begin{simplebox}
        \small
        \noindent
        \begin{minipage}[t]{0.27\textwidth}
            \begin{algorithmic}[1]
                \Statex \textbf{$\Gen(1^n)$}
                \Indent
                \State \Return $\Gen^0(1^n)$
                \EndIndent
            \end{algorithmic}
        \end{minipage}
        \hfill
        \begin{minipage}[t]{0.30\textwidth}
            \begin{algorithmic}[1]
                \Statex \textbf{$\Enc_\pk(k ,m)$}
                \Indent
                    \State \Return $\left(\Enc^0_\pk(k,m_n^0), \Enc^0_\pk(k_n^0 ,m)\right)$
                \EndIndent
            \end{algorithmic}
        \end{minipage}
        \hfill
        \begin{minipage}[t]{0.40\textwidth}
            \begin{algorithmic}[1]
                \Statex \textbf{$\Dec_\sk(ct_1, ct_2)$}
                \Indent
                    \If{$\Dec^0_\sk(ct_1)=\bot\lor \Dec^0_\sk(ct_2)=\bot$}
                        \State \Return $\bot$
                    \EndIf
                    \State $(k_1, m_1)\coloneqq \Dec^0_\sk(ct_1)$
                    \State $(k_2, m_2)\coloneqq \Dec^0_\sk(ct_2)$
                    \If{$m_1\neq m^0_n \lor k_2\neq k^0_n$}
                        \State \Return $\bot$
                    \EndIf
                    \State \Return $(k_1, m_2)$
                \EndIndent
            \end{algorithmic}
        \end{minipage}
    \end{simplebox}
     We claim that $\Pi=(\Gen, \Enc, \Dec)$ is indeed correct and IND-CPA-secure.
    \begin{claim}\label{claim:p1-1}
        $\Pi=(\Gen, \Enc, \Dec)$ is a correct PKE scheme. Specifically, $\Pi$ is $2\delta$-correct. 
    \end{claim}
    \begin{proof}
        Since the PKE scheme $\Pi^0=(\Gen^0, \Enc^0, \Dec^0)$ is $\delta$-correct, for each $n\in\Nbb$ and any $(k,m)\in\Kcal_n\times \Mcal_n$,
        \begin{equation}\label{eq:p1-1}
            \Pr\left[\Dec^0_\sk(\Enc^0_\pk(k,m))\neq (k,m) \right] \le \delta,
        \end{equation}
        where the probability is taken over the randomness of $\Enc^0_\pk$ and $(\pk,\sk)\sample\Gen^0(1^n)$.

        Therefore, for each $n\in\Nbb$ and any $(k,m)\in\Kcal_n\times \Mcal_n$, we deduce by union bound that
        \begin{align*}
            \Pr\left[\Dec_\sk(\Enc_\pk(k,m))\neq (k,m) \right]
            &= \Pr\left[\Dec^0_\sk(\Enc^0_\pk(k,m^0_n))\neq (k,m^0_n) \lor \Dec^0_\sk(\Enc^0_\pk(k^0_n, m))\neq (k^0_n, m)  \right]\\
            &\le \underbrace{\Pr\left[\Dec^0_\sk(\Enc^0_\pk(k,m^0_n))\neq (k,m^0_n)\right]}_{\begin{gathered}\le \delta \;\; \textrm{(by \eqref{eq:p1-1})}\end{gathered}}+\underbrace{\Pr\left[\Dec^0_\sk(\Enc^0_\pk(k^0_n, m))\neq (k^0_n, m)  \right]}_{\begin{gathered}\le \delta\;\; \textrm{(by \eqref{eq:p1-1})}\end{gathered}}\\
            &\le 2\delta,
        \end{align*}
        which is negligible as $\delta$ is negligible. Hence, by definition, $\Pi$ is correct and also $2\delta$-correct. 
    \end{proof}

    \begin{claim}\label{claim:p1-2}
        The PKE scheme $\Pi=(\Gen, \Enc, \Dec)$ is IND-CPA-secure.
    \end{claim}
    \begin{proof}
        We prove by contradiction. Suppose to the contrary that  $\Pi=(\Gen, \Enc, \Dec)$ is not IND-CPA-secure. This means there exist a PPT adversary $\Adv$ and a nonnegligible function $\epsilon(\cdot)$ such that in the following experiment $\PKE^{\textsf{IND-CPA}}_{\Pi}$, $\Adv$ wins with advantage $\epsilon(n)$, i.e., 
        \begin{equation}\label{eq:p1-2}
            \Pr[\Adv \textrm{ wins } \PKE^{\textsf{IND-CPA}}_{\Pi}] 
            \equiv \Pr_{\substack{((k_0, m_0), (k_1, m_1))\sample \Adv(\pk)\\b\Usample\bit}}\left[\Adv\left(\Enc^0_\pk(k_b, m_n^0),\Enc^0_\pk(k_n^0, m_b)\right) = b\right]
            = \dfrac{1}{2}+\epsilon(n),
        \end{equation}
        where the probability is taken over $(\pk, \sk)\sample\Gen^0(1^n)$.
        \end{proof}\begin{proof}[cont'd]
        {
            \[
                \begin{array}{@{} c c c @{}}
                \Ch & & \Adv \\
                (\pk, \sk)\sample \Gen(1^n)=\Gen^0(1^n) & \XR{\pk}  &\\
                b\Usample \bit& \XL{((k_0, m_0), (k_1, m_1))} &\\
                (ct_1, ct_2)\coloneqq \Enc_\pk(k_b, m_b) = \left(\Enc^0_\pk(k_b, m_n^0),\Enc^0_\pk(k_n^0, m_b)\right) & \XR{(ct_1, ct_2)} &\\
                \textrm{$\Adv$ wins if $b'=b$} & \XL{b'} & \\
                \end{array}
            \]
            \vspace{-0.8cm}
            \captionof{figure}{Experiment $\PKE^{\textsf{IND-CPA}}_{\Pi}$}
            \label{fig:p1-game1}
        }

        We construct a PPT adversary (a reduction) $R$ that breaks the IND-CPA security of $\Pi^0=(\Gen^0, \Enc^0, \Dec^0)$. The construction of $R$ is provided in the following experiment $\PKE^{\textsf{IND-CPA}}_{\Pi^0}$. Worth noting, $R$ uniformly selects $i\in\bit$ and behaves differently according to $i$ because this allows us to analyze the advantage of $R$ by hybrid argument.
        {
            \[
                \begin{array}{@{} c c c c c @{}}
                \Ch & & R & & \Adv \\
                (\pk, \sk)\sample \Gen^0(1^n) & \XR{\pk} & i\Usample \bit & \XR{\pk} &\\
                b\Usample \bit& \XL{((k^\bigstar_0, m^\bigstar_0), (k^\bigstar_1, m^\bigstar_1))} &
                \begin{aligned}
                &(k^\bigstar_0, m^\bigstar_0), (k^\bigstar_1, m^\bigstar_1) \\
                &\qquad\coloneqq {\begin{cases}
                    (k_0, m^0_n), (k_1, m^0_n) & \textrm{if } i=0\\
                    (k^0_n, m_0), (k^0_n, m_1) & \textrm{if } i=1\\
                \end{cases}} \end{aligned} & \XL{((k_0, m_0), (k_1, m_1))} &\\
                ct\coloneqq \Enc^0_\pk(k^\bigstar_b, m^\bigstar_b) & \XR{ct} & (ct_1, ct_2)\coloneqq \begin{cases}
                    \left(ct, \Enc^0_\pk(k^0_n, m_1)\right) & \textrm{if } i=0\\
                    \left(\Enc^0_\pk(k_0, m^0_n), ct\right) & \textrm{if } i=1
                \end{cases} & \XR{(ct_1, ct_2)}&\\
                \textrm{$R$ wins if $b'=b$} & \XL{b'} & & \XL{b'}&\\
                \end{array}
            \]
            \vspace{-0.8cm}
            \captionof{figure}{Experiment $\PKE^{\textsf{IND-CPA}}_{\Pi^0}$}
            \label{fig:p1-game2}
        }

        Clearly $R$ runs in PPT as $\Adv$ is a PPT algorithm. In the experiment $\PKE^{\textsf{IND-CPA}}_{\Pi^0}$, for all $n\in\Nbb$ the win rate of $R$ is
        \begin{align*}
            &\Pr\left[R \textrm{ wins } \PKE^{\textsf{IND-CPA}}_{\Pi^0}\right]
            = \Pr_{i\Usample\bit, b\Usample\bit}\left[\Adv\left(ct_1, ct_2\right) = b\right]\\
            =&\frac{1}{2}\Pr_{b\Usample\bit}\left[\Adv\left(\Enc^0_\pk(k_b, m^0_n), \Enc^0_\pk(k^0_n, m_1) \right) = b\right] + \frac{1}{2} \Pr_{b\Usample\bit}\left[\Adv\left(\Enc^0_\pk(k_0, m^0_n), \Enc^0_\pk(k^0_n, m_b) \right) = b\right]\\
            =&\frac{1}{2}\left(\frac{1}{2} \Pr\left[\Adv\left(\Enc^0_\pk(k_1, m^0_n), \Enc^0_\pk(k^0_n, m_1) \right) = 1\right] 
            + \frac{1}{2} \highlight{CarnationPink}{\Pr\left[\Adv\left(\Enc^0_\pk(k_0, m^0_n), \Enc^0_\pk(k^0_n, m_1) \right) = 0\right]} \right)\\
            +& \frac{1}{2}\left(\frac{1}{2} \highlight{CarnationPink}{\Pr\left[\Adv\left(\Enc^0_\pk(k_0, m^0_n), \Enc^0_\pk(k^0_n, m_1) \right) = 1\right]}
            +\frac{1}{2} \Pr\left[\Adv\left(\Enc^0_\pk(k_0, m^0_n), \Enc^0_\pk(k^0_n, m_0) \right) = 0\right]\right)\\
            =& \frac{1}{4}\cdot \highlight{CarnationPink}{1}+\frac{1}{2}\left(\frac{1}{2}\Pr\left[\Adv\left(\Enc^0_\pk(k_0, m^0_n), \Enc^0_\pk(k^0_n, m_0) \right) = 0\right]+\frac{1}{2}\Pr\left[\Adv\left(\Enc^0_\pk(k_1, m^0_n), \Enc^0_\pk(k^0_n, m_1) \right) = 1\right]\right)\\
            =& \frac{1}{4}+\frac{1}{2}\Pr_{b\Usample\bit}\left[\Adv\left(\Enc^0_\pk(k_b, m^0_n), \Enc^0_\pk(k^0_n, m_b) \right) = b\right]\\
            =& \frac{1}{4}+\frac{1}{2}\left(\frac{1}{2}+\epsilon(n)\right) \qquad \textrm{(by \eqref{eq:p1-2})}\\
            =&\frac{1}{2}+\frac{1}{2}\epsilon(n),
        \end{align*}
        where all probabilities are taken over $(\pk, \sk)\sample\Gen^0(1^n),\;\; ((k_0, m_0), (k_1, m_1))\sample \Adv(\pk)$. 
        
        Since the advantage $\frac{1}{2}\epsilon(n)$ of the PPT adversary $R$ in the experiment $\PKE^{\textsf{IND-CPA}}_{\Pi^0}$ is nonnegligible, by definition,  $\Pi^0$ is not IND-CPA-secure, which contradicts our assumption.
    \end{proof}

    Now it suffices to show that the corresponding $\Pi'=(\Gen', \Enc', \Dec')$ is not IND-CCA-secure. To show this, we construct an oracle adversary $\Adv^{\Pi'}$ that breaks the IND-CCA security of $\Pi'=(\Gen', \Enc', \Dec')$. The construction of $\Adv^{\Pi'}$ is given in the following experiment $\PKE^{\textsf{IND-CCA}}_{\Pi'}$, where we have expanded the definitions of $\Gen', \Enc', \Dec'$ for convenience.
    \begin{idea}
        Our main idea is to replace the encrypted key $k$ in the received ciphertext $ct$ with another key $k'$, sign the new, distinct ciphertext $ct'$ with key $k'$ to obtain a valid tag $t'$, and then query $(ct', t')$ to the decryption oracle $\Dec'_\sk(\cdot)$, where $(ct', t')\neq (ct, t)$.
    \end{idea}
    {
            \[
                \begin{array}{@{} c c c @{}}
                \Ch & & \Adv^{\Pi'} \\
                (\pk, \sk)\sample \Gen(1^n)=\Gen^0(1^n) & \XR{\pk}  &\\
                & \XL{(m_0, m_1)} & \makecell[c]{\textrm{Sample } m_0,m_1\Usample \Mcal_n \\\textrm{ where } m_0\neq m_1}\\
                b\Usample \bit &&\\
                k\Usample\Kcal_n &&\\
                ct=(ct_1, ct_2)\sample \Enc_\pk(k, m_b) = \left(\Enc^0_\pk(k, m^0_n), \Enc^0_\pk(k^0_n, m_b) \right) &&\\
                t\sample \Sign_k(ct) &  & \\
                & \XR{(ct=(ct_1, ct_2), t)} &\\
                && k'\Usample \Kcal_n\\
                && ct' \sample \left(\Enc^0_\pk(k', m^0_n), ct_2 \right)\\
                && t' \sample \Sign_{k'}(ct')\\
                && \textrm{If } (ct', t')=(ct, t) \textrm{ Then ABORT}\\
                \makecell[c]{
                    (k'', m')\coloneqq \Dec_\sk(ct')\\
                    \tilde{m}\coloneqq {\begin{cases}
                        m' &\textrm{if } (k'', m')\neq \bot \land \Vrfy_{k''}(ct', t')=\accept\\
                        \bot &\textrm{otherwise}
                    \end{cases}}
                } & \makecell[c]{
                    \XL{(ct', t')\neq (ct, t)}\\
                    \\\\
                    \XR{\tilde{m}}
                } & \textrm{Query } \Dec'_\sk(\cdot)\\
                && b'\coloneqq \begin{cases}
                    0 & \textrm{if } \tilde{m} = m_0\\
                    1 & \textrm{if } \tilde{m} = m_1\\
                    \Usample \bit & \textrm{otherwise} 
                \end{cases}\\
                \textrm{$\Adv^{\Pi'}$ wins if $b'=b$} & \XL{b'} & \\
                \end{array}
            \]
            \captionof{figure}{Experiment $\PKE^{\textsf{IND-CCA}}_{\Pi'}$}
            \label{fig:p1-game3}
        }

        Clearly $\Adv^{\Pi'}$ is a PPT algorithm.
        
        Recall that $\Pi^0=(\Gen^0, \Enc^0,\Dec^0)$ is $\gamma$-spread where $2^{-\gamma}=\negl(n)$, meaning that for each $n\in\Nbb$ and any $(k,m)\in \Kcal_n \times \Mcal_n, ct\in\Ccal_n$, we have
        \begin{equation}\label{eq:p1-3}
            \Pr\left[ct=\Enc^0_\pk(k,m)\right] \le 2^{-\gamma} = \negl(n),
        \end{equation}
        where the probability is taken over the randomness of $\Enc^0_\pk$ and $(\pk, \sk)\sample \Gen^0(1^n)$. Using this, we show that the probability that $\Adv^{\Pi'}$ aborts, denoted $\Pr[\textrm{ABORT}]$, is negligible:
        \begin{equation}\label{eq:p1-5}
            \Pr[\textrm{ABORT}]
            = \Pr[(ct'=t')=(ct,t)]
            \le \Pr[ct'=ct]
            \le \Pr\left[\Enc^0_\pk(k', m^0_n)=ct_1\right]
            \stackrel{\eqref{eq:p1-3}}\le 2^{-\gamma}
            =\negl(n).
        \end{equation}

        Moreover,  since $\Pi=(\Gen, \Enc,\Dec)$ is $2\delta$-correct (which is proven by \autoref{claim:p1-1}) and  $\delta=\negl(n)$, meaning that for each $n\in\Nbb$ and $(k,m)\in\Kcal_n\times \Mcal_n$
        \begin{equation}\label{eq:p1-4}
            \Pr\left[\Dec_\sk\left(\Enc_\pk(k,m)\right)\neq (k,m)\right] \le 2\delta = \negl(n),
        \end{equation}
        we see that the probability that $ct'$ is correctly decrypted back to $(k', m_b)$ is $1-\negl(n)$:
        \begin{align}\label{eq:p1-6}
            \Pr[(k'', m')=(k', m_b)]
            &= \Pr\left[\Dec_\sk\left(\Enc^0_\pk(k', m^0_n), \Enc^0_\pk(k^0_n, m_b)\right)=(k', m_b)\right]\notag\\
            &= \Pr\left[\Dec_\sk\left(\Enc_\pk(k', m_b)\right)=(k', m_b)\right]\notag\\
            &\stackrel{\eqref{eq:p1-4}}\ge 1-2\delta = 1-\negl(n).
        \end{align}

        Also, conditioning on $(k'', m')=(k', m_b)$ (actually $k''=k'$ suffices), the probability that $\Vrfy_{k''}(ct',t')=\accept$ (i.e. MAC tag $t'$ is valid for message $ct'$) is $1$ by the correctness of the MAC system $\Ical=(\Sign, \Vrfy)$:
        \begin{equation}\label{eq:p1-7}
            \Pr\left[\Vrfy_{k''}(ct',t')=\accept \mid (k'', m')=(k', m_b) \right]
            = \Pr\left[\Vrfy_{k'}(ct',\Sign_{k'}(ct'))=\accept \right] =1.
        \end{equation}
    
        Putting it all together, we see that in the experiment  $\PKE^{\textsf{IND-CCA}}_{\Pi'}$,  since $m_0\neq m_1$, $\Adv^{\Pi'}$ wins if $\tilde{m}=m_b \land \neg \textrm{ABORT}$, i.e., $\Adv^{\Pi'}$ correctly obtains $m_b$ via the decryption oracle without aborting. Therefore, for all $n\in\Nbb$
        \begin{align*}
            \Pr[\Adv^{\Pi'} \textrm{ wins } \PKE^{\textsf{IND-CCA}}_{\Pi'}]
            &\ge \Pr\left[ \tilde{m}=m_b \land \neg \textrm{ABORT}\right]\\
            &= \Pr\left[ m'=m_b\land (k'', m')\neq \bot \land \Vrfy_{k''}(ct',t')=\accept  \land \neg\textrm{ABORT}\right]\\
            &\ge  \Pr\left[ (k'', m')=(k', m_b) \land \Vrfy_{k''}(ct',t')=\accept  \land \neg\textrm{ABORT}\right]\\
            &\ge \Pr\left[ (k'', m')=(k', m_b) \land \Vrfy_{k''}(ct',t')=\accept \right] - \Pr[\textrm{ABORT}]
            \qquad (\textrm{by Union Bound})\\
            &= \Pr\left[ (k'', m')=(k', m_b)\right] \Pr\left[\Vrfy_{k''}(ct',t')=\accept \mid  (k'', m')=(k', m_b)\right] - \Pr[\textrm{ABORT}]\\
            &\ge (1-\negl(n))\cdot 1 - \negl(n)
            \qquad \textrm{(by \eqref{eq:p1-5}, \eqref{eq:p1-6}, \eqref{eq:p1-7})}\\
            &= \frac{1}{2}+\nonnegl(n).
        \end{align*}
        This implies that the advantage of $\Adv^{\Pi'}$ in the experiment  $\PKE^{\textsf{IND-CCA}}_{\Pi'}$ is nonnegligible.

        Hence, by definition, $\Pi'=(\Gen', \Enc', \Dec')$ is not IND-CCA-secure.
\end{answer}